\documentclass[upright, contnum]{umemoria}

%fix for the oneside argument
\makeatletter
\g@addto@macro\titlepage{\pagenumbering{Alph}}
\g@addto@macro\endtitlepage{\pagenumbering{roman}}
\makeatother

\depto{Departamento de Ingeniería Industrial}
\author{Andrea Ignacia Canales Gutierrez}
\title{Essays on Anti-Competitive Behavior}
\auspicio{}
\date{Octubre 2018}
\guia{Juan Escobar Castro}
\carrera{Doctora en Sistemas de Ingeniería}
\memoria{Tesis para optar al Grado de \break  Doctora en Sistemas de Ingeniería}
\comision{Leonardo Basso Sotz, Gastón Llanes, Matteo Triossi Verondini}

\usepackage{lipsum}

\usepackage[utf8]{inputenc}
\usepackage[T1]{fontenc}

\begin{document}

\frontmatter
\maketitle

\begin{abstract}
{En esta tesis, a raíz de dos problemas distintos, se estudia cómo la entrada de nuevas firmas a un mercado puede modificar los incentivos a desarrollar conductas anticompetitivas por parte de la o las firmas incumbentes. En primer lugar, estudiamos un problema de colusión con entrada de nuevas firmas. Por otra parte, se modelan los incentivos a utilizar estrategias predatorias cuando los costos de producción pueden disminuir o aumentar, dependiendo si logran acumular experiencia en el mercado o no.
En el capítulo 1 se estudian los carteles incompletos. La evidencia muestra que, en muchos casos, los acuerdos colusivos no se conforman con la participación de todas las firmas del mercado si no que, un grupo de firmas convive con una franja competitiva de empresas, estos se denominan carteles incompletos. En este contexto estudiamos el acuerdo colusivo óptimo, donde las variaciones de precio, por parte de las firmas que lo fijan, pueden generar entrada o salida de firmas de la franja, y por tanto, cambiar la estructura del mercado. En este modelo, una firma que se desvía del acuerdo colusivo puede tener incentivos a fijar un precio muy bajo, con el objetivo de sacar del mercado firmas de la franja competitiva. Esta salida se produce solo si las firmas de la franja esperan un castigo severo por parte del cartel ante el desvío, es decir,una etapa de castigo donde se observarían precios muy bajos. Este trabajo identifica un nuevo trade-off entre la severidad del castigo y las ganancias del desvío de un acuerdo colusivo.
En el capítulo 2 se aborda un problema de predación. Se ha documentado que en al- gunos mercados, a medida que las firmas acumulan experiencia en la producción de bienes y/o servicios, disminuyen sus costos, este fenómeno de denomina learning-by-doing. Por otro lado, también se documenta que la pérdida de experiencia produce el fenómeno contrario, llamado organizational forgetting. En este capítulo de desarrolla un modelo general en un mercado donde pueden coexistir ambos efectos y se estudian los incentivos a utilizar estra- tegias predatorias. Se considera un modelo asimétrico, donde la firma incumbente tiene un mercado cautivo, y puede forzar la salida de una firma entrante. Se observa que ambas firmas pueden utilizar estrategias agresivas de precios pero con distintos objetivos: la incumbente para forzar la salida de la entrante, y esta última para evitar ser predada. Además, para un caso reducido, observamos que la interacción de learning-by-doing y organizational forgetting provoca una competencia en precios menos agresiva que si solo está presente uno de estos efectos en el mercado.}
\end{abstract}

\begin{dedicatoria} % opcional
Una dedicatoria corta. Por ejemplo, al \emph{Centro Tecnológico Ucampus}
\end{dedicatoria}

\begin{thanks} % opcional
\lipsum[1-2]
\end{thanks}
\cleardoublepage

\tableofcontents
\listoftables % opcional
\listoffigures % opcional

\mainmatter

\chapter{Introduction}

% En las últimas décadas, el uso de la ciencia de datos ha transformado la forma en que las organizaciones toman decisiones. Las compañías más exitosas destacan por su capacidad para tomar decisiones informadas y motivadas por el análisis de datos, lo que les permite anticiparse a cambios, optimizar procesos y mantenerse competitivas en el mercado. Del mismo modo, los gobiernos e instituciones públicas han comenzado a aprovechar estas herramientas, adoptando un enfoque de toma de decisiones basado en datos. Este cambio permite a las instituciones enfrentar problemas complejos con mayor precisión y efectividad, desde la gestión de crisis sanitarias hasta la mejora de procesos en la administración pública.

% En este contexto, la presente tesis se centra en dos áreas aplicadas de vital importancia para el desarrollo social y económico de los países: la gestión de la pandemia de COVID-19 y el rediseño de mecanismos de contratación pública. Estos dos casos muestran cómo el uso de datos no solo permite mejorar la toma de decisiones en situaciones críticas, sino también aumentar la eficiencia, transparencia y equidad en los procesos del Estado.

% La primera investigación se centra en un análisis detallado de la movilidad durante la pandemia en la ciudad de Santiago, lo que permitió identificar disparidades en la respuesta a las políticas de confinamiento y distanciamiento social. (Antes: En primer lugar, el análisis detallado de la movilidad durante la pandemia en la ciudad de Santiago permitió identificar disparidades en la respuesta a las políticas de confinamiento y distanciamiento social.) Se observó que la adherencia a las medidas sanitarias variaba significativamente según el nivel socioeconómico de las comunas. Este análisis granular fue clave para entender por qué ciertos sectores experimentaron mayores tasas de infección, lo que, a su vez, proporcionó información crucial para diseñar políticas de salud pública más focalizadas y efectivas.

% En segundo lugar, la tesis aborda la necesidad de optimizar los procesos de contratación pública, con especial atención al ``Convenio Marco'' en Chile. Utilizando técnicas avanzadas de ciencia de datos, se identificaron ineficiencias y cuellos de botella en el sistema, lo que permitió proponer reformas que mejoran la transparencia y la eficiencia del gasto público. Esta área de estudio es de especial relevancia, ya que un uso más eficiente de los recursos públicos puede traducirse en mejores servicios para la población.

% Ambos casos reflejan cómo la ciencia de datos puede ser una herramienta poderosa para resolver problemas estructurales en el sector público, mejorando la eficiencia en tiempos de crisis y en el funcionamiento cotidiano del Estado. Esta investigación destaca la importancia de desarrollar capacidades analíticas en las instituciones para enfrentar desafíos futuros de manera más informada y eficiente.

% (EN - Ingles)

In recent decades, the use of data science has transformed the way organizations make decisions. The most successful companies excel in their ability to make informed, data-driven decisions, allowing them to anticipate changes, optimize processes, and remain competitive in the market. Similarly, governments and public institutions have started using these tools, adopting a data-driven approach to decision-making. This shift enables institutions to tackle complex problems with greater precision and effectiveness, ranging from managing public health crises to improving administrative processes.

In this context, this thesis focuses on two applied areas of vital importance for the social and economic development of countries: the management of the COVID-19 pandemic and the redesign of public procurement mechanisms. These two cases demonstrate how the use of data not only improves decision-making in critical situations but also enhances efficiency, transparency, and fairness in government processes.

The first study focuses on a detailed analysis of mobility during the pandemic in the city of Santiago, which allowed for the identification of disparities in the response to lockdown and social distancing policie. It was observed that adherence to public health measures varied significantly based on the socioeconomic status of different neighborhoods. This granular analysis was key to understanding why certain areas experienced higher infection rates, which in turn provided crucial insights for designing more targeted and effective public health policies.

The second study addresses the need to optimize public procurement processes, with a special focus on the “Convenio Marco” in Chile. Using advanced data science techniques, inefficiencies and bottlenecks in the system were identified, leading to proposals for reforms that improve transparency and the efficiency of public spending. This area of study is particularly relevant, as more efficient use of public resources can translate into better services for the population.

Both cases reflect how data science can be a powerful tool to solve structural problems in the public sector, improving efficiency in times of crisis and in the everyday functioning of the state. This research highlights the importance of developing analytical capacities within institutions to better address future challenges in a more informed and efficient manner.
\chapter{Primero}
\lipsum[1-3]
\begin{defn}[ver \cite{KAR00}] Definición definitiva $$\frac{d}{dx}\int_a^xf(y)dy=f(x).$$\end{defn}
\chapter{Segundo}
\lipsum[50-60]
\begin{conclusion}
	\lipsum[130-132]
	\begin{figure}[!h]
		\centering
		\includegraphics[scale=.2]{imagenes/fcfm.pdf}
		\caption{Logo de la Facultad}
		\label{logofcfm}
	\end{figure}
	\lipsum[133-134]
	\begin{table}[!h]
		\centering
		\begin{tabular}{|c||c|}
			\hline
			Campo 1& Campo 2\\\hline
			Valor 1& Valor2\\\hline
		\end{tabular}
		\caption{Tabla 1}
		\label{tabla:1}
	\end{table}
	\lipsum[135]
\end{conclusion}


% \input{glosario.tex} % opcional

\bibliographystyle{plain}
\bibliography{bibliografia}

% % \chapter*{Annexes}
% \addcontentsline{toc}{chapter}{Annexes}

% \chapter{Further details of methodologies, equations and robustness analysis}


\chapter{Econometric Model and Data Processing} \label{app:data}


\section{Model specification}

This section provides a detailed description of the econometric models that were estimated using a difference-in-difference design.

Let the index $t\in\{\text{Old},\text{New}\}$ represent the old (2014) and the new (2017) Food FAs, respectively.  Let $b_{ijrt}$ represent the bid offered in FA $t$ by supplier $j$ to provide product $i$ in region $r$, including shipping.  Define $B^{med}_{irt}$ as the \textit{median} bid across all supplier bids submitted in the FA for that product-region-FA combination. In calculating this median bid price, we considered alternative methods for discarding extreme prices that could be generated by bidding mistakes or unrealistically aggressive bidding. For the main results we used a modified Tukey rule \citep{Tukey1977}, described in the \textit{Data collection and processing} section. To make the bids for the 2014 and 2017 Food FAs comparable, all prices for 2014 were adjusted using the CPI food price index. In addition, we normalized bid prices so that they represented bid prices per unit. See the \textit{Data collection and processing} section for further details on data pre-processing. 

The following difference-in-differences specification is used to estimate the effect of the competitive treatment condition ---denoted by the indicator variable $Comp_{i}$--- on bid prices:
\begin{equation}
    \log (B^{med}_{irt}) = \delta_r + \gamma_i + \alpha New_{t} + \beta New_{t}\times Comp_{i} + \varepsilon_{irt} \ ,
    \label{eq:reg_sub_bid}
\end{equation}

\noindent where $New_t$ is an indicator equal to one for the new FA and $\delta_r$ is a region fixed effect. The sample is comprised of all products $i$ that were matched across the old and new FAs based on similar attributes, and the regression includes a product fixed effect $\gamma_i$. {The error term $\varepsilon_{irt}$ represents other idiosyncratic unobservable factors that affect bid prices of each product across FAs and  regions.} The coefficient $\alpha$ measures the average differences in bid prices between the old and new FAs for the noncompetitive baseline group. The key parameter of interest is $\beta$, the coefficient capturing the incremental change in bid prices for the auctions in the competitive treatment condition in the new FA.

The second model specification uses panel data on posted and transaction prices. Let $p_{ijrtw}$ denote the average posted prices by supplier $j$ for product $i$ in region $r$ during calendar week $w$ in FA $t$. As with the bids, we calculated the {\em median posted price} for a product across all suppliers during that week, denoted by $P^{med}_{irtw}$. That is, $P^{med}_{irtw}$ denotes the median price at which product $i$ could be purchased in region $r$ during calendar week $w$ under the operation of FA $t$ in the period listed above. 

We estimated the panel regression:
\begin{equation}
    \log (P^{med}_{irtw})= \delta_r + \gamma_i + \tau_{wc(i)} + \alpha New_{t} + \beta New_{t}\times Comp_{i} + \varepsilon_{irtw} \ ,
    \label{eq:reg_op_posted}
\end{equation}
\noindent where $\delta_r$ and $\gamma_i$ are region and product fixed effects respectively, $\tau_{wc(i)}$ is a product category-specific calendar week dummy variable to capture potential seasonality ($c(i)$ indicates the category of product $i$). {The error term $\varepsilon_{irtw}$ represents other unobservable factors that affect prices during the operation stage.} All the models were estimated using Ordinary Least Squares and robust standard errors. 

\section{Data collection and processing}
In order to compare prices between different FA and to rule out unreasonable values, we establish a procedure to process bids, transaction prices, and posted prices.

\paragraph{Bids.}  In the case of bids, we identify outliers within each distribution of bids defined by a standardized product $i$ and FA $t$. In the case of FA 2014, we use prices before adding the respective shipping rate. In the case of FA 2017, we normalize prices to the Metropolitan Region (where the capital city Santiago is located) for outlier identification. Additionally, we use the CPI index to adjust prices from FA 2014 for inflation. In Chile there are available CPI indexes that are product specific, which we use to adjust the prices offered in 2014 to present value (to the year 2017), and therefore each product category is adjusted differently according to the corresponding specific CPI. 

We identify four types of outliers, using four different approaches. For each FA $t$ and standardized product $i$, we define a set of similar products $C_{i,t}$. (A standardized product was defined in the main text.) First, we compute the mean price $\hat{p}_{i,t}$ of its distribution. For each supplier $s$, product $j$, region $r$, and FA $t$, with bid $p_{s,j,r,t}$, we compute the ratio $R_{s,j,r,t} = \frac{p_{s,j,r,t}}{\hat{p}_{i,t}}$, where $i$ is the standardized product associated with $j$. We define outliers as follows:

\begin{itemize}
\item[] \textbf{Type 1}: $p_{s,j,r,t}$ is an outlier if $R_{s,j,r,t}<1/2$ or $R_{s,j,r,t}>2$.
\item[] \textbf{Type 2}: $p_{s,j,r,t}$ is an outlier if $R_{s,j,r,t}<1/3$ or $R_{s,j,r,t}>3$.
\item[] \textbf{Type 3}: $p_{s,j,r,t}$ is identified as an outlier by Tukey's rule: observations that more than 1.5$\times$IQR (inter-quartile range) below and above the first and third quartile, respectively.
\item[] \textbf{Type 4}: If 20\% of the products in the set of similar products $C_{i,t}$ meet the Type 3 outlier criterion, all products in the set are considered outliers. That is, we exclude all the standardized products that exhibit a large proportion (20\% or more) of extreme values.%\todoMO{Re-wrote explanation}
\end{itemize}

First, to compute normalized bids or prices, we divide prices by the relevant number of units in the bundle.
In addition, as a robustness check we estimate regression \eqref{eq:reg_sub_bid} by correcting prices by volume to take into account potential scale effects such as volume discounts. See Tables \ref{tab:estsub_vol} and \ref{tab:estadj_vol}. Thus, prices of products specifying larger quantities of items would be comparable to those of products with a single or a few items. Specifically, in FA 2014 products were sold in large bundles. To account for potential volume discounts, we estimate 
 the regression: 
\begin{equation}
    \log (B^{med}_{irt}) = \delta_r + \gamma_i + \beta \log(units_{irt}) + \varepsilon_{irt} \ ,
    \label{eq:reg_sub_bid_rob}
\end{equation}
where $units_{irt}$ represents the amount of items per bundle for product $i$. Then, we compute the adjusted bid price  $\hat{B}^{med}_{irt} = \log B^{med}_{irt} - \hat{\beta}\log(units_{irt})$ and use this as the dependent variable in regression \eqref{eq:reg_sub_bid}.

\paragraph{Transaction and posted prices.} Regarding transaction and posted prices, we correct by inflation according to the monthly CPI index for Food and Non-alcoholic Beverages, using as the reference month December 2013. For transaction and posted prices, we identify three types of outliers. Let be $C(i,t)$ the set of transaction/posted prices of a standardized product $i$ in FA $t$. For each transaction/post $k$ we define $R_{k,t} = \frac{p_{k,t}}{\hat{p}_{i,t}}$, where $p_{k,t}$ is the unit price of $k$, and $\hat{p}_{i,t}$ is the average price in $C_{i,t}$. Then, we define outliers as follows:

\begin{itemize}
\item[] \textbf{Type 1}: If the proportion of transactions/posts with $R_{k,t}<1/2$ or $R_{k,t}>2$ is greater than 20\%.
\item[] \textbf{Type 2}: If the proportion of transactions/posts with $R_{k,t}<1/3$ or $R_{k,t}>3$ is greater than 20\%.
\item[] \textbf{Type 3}: If the proportion of outliers in $C_{i,t}$ is greater than 20\% by Tukey's rule.
\end{itemize}

As a robustness check, in the \textit{Robustness analysis and alternative regression specifications} section we show the results under the different outlier elimination methods.


\chapter{Robustness analysis and alternative regression 
 specifications}\label{app:robustness}

%- Results with different forms of outliers
%- Results with TP clustering
%- Results with log(q) adjustments for bids.

\section{Regressions for submitted bids}



\begin{table}[H]
   \centering
   \begin{tabular}{lcccc}
      \toprule
                                & Type 1         & Type 2         & Type 3         & Type 4 \\   
                                & (1)            & (2)            & (3)            & (4)\\  
      \midrule 
      New                       & -0.148$^{***}$ & -0.151$^{***}$ & -0.146$^{***}$ & -0.141$^{***}$\\   
                                & (0.005)        & (0.005)        & (0.005)        & (0.006)\\   
      New $\times$ Comp  & -0.029$^{***}$ & -0.034$^{***}$ & 0.009          & 0.002\\   
                                & (0.007)        & (0.008)        & (0.009)        & (0.010)\\   
       \\
      Observations              & 13,111         & 13,195         & 13,163         & 12,349\\  
      R$^2$                     & 0.94644        & 0.93650        & 0.92588        & 0.92289\\  
      Adjusted R$^2$            & 0.94433        & 0.93401        & 0.92297        & 0.91965\\  
      \bottomrule
   \end{tabular}
   
   \par \raggedright 
   \caption{Estimation for submitted bids under different outlier elimination rules. This table shows estimates for submitted bids according to  \eqref{eq:reg_sub_bid}. Each model is defined by the type of method used for outlier detection according to section \ref{app:data}. Standardized product and region fixed effect are included in all Models. We compute \textit{robust} standard errors.}
\end{table}




\begin{table}[H]
   \label{app:tab:volume_adj}
   \bigskip
   \centering
   \begin{tabular}{lcccc}
      \toprule
                                & Type 1         & Type 2         & Type 3         & Type 4 \\   
                                & (1)            & (2)            & (3)            & (4)\\  
      \midrule 
      New                       & -0.125$^{***}$ & -0.128$^{***}$ & -0.176$^{***}$ & -0.177$^{***}$\\   
                                & (0.005)        & (0.005)        & (0.005)        & (0.005)\\   
      New $\times$ Comp  & -0.023$^{***}$ & -0.028$^{***}$ & 0.001          & -0.008\\   
                                & (0.007)        & (0.008)        & (0.009)        & (0.009)\\   
       \\
      Observations             & 13,111         & 13,195         & 13,163         & 12,349\\  
      R$^2$                     & 0.94616        & 0.93629        & 0.92623        & 0.92344\\  
      Adjusted R$^2$            & 0.94403        & 0.93379        & 0.92333        & 0.92023\\  
      \bottomrule
   \end{tabular}
   
   \par \raggedright 
   \caption{\label{tab:estsub_vol} Estimation for submitted bids: Adjusting by volume. This table shows estimates for submitted bids  when prices are adjusted by volume following the specification in \eqref{eq:reg_sub_bid_rob}. Each model is defined by the type of method used for outlier detection according defined in Appendix~\ref{app:data}. Standardized product and region fixed effect are included in all models. We compute robust standard errors.}
\end{table}





\begin{table}[H]
   \centering
   \begin{tabular}{lcccc}
      \toprule
       & Type 1         & Type 2         & Type 3         & Type 4 \\ & (1)            & (2)            & (3)            & (4)\\  
      \midrule 
      New                       & -0.123$^{***}$ & -0.125$^{***}$ & -0.118$^{***}$ & -0.115$^{***}$\\   
                                & (0.005)        & (0.005)        & (0.006)        & (0.006)\\   
      New $\times$ Comp  & -0.041$^{***}$ & -0.048$^{***}$ & -0.001         & -0.012\\   
                                & (0.008)        & (0.009)        & (0.010)        & (0.010)\\   
       \\
      Observations              & 11,760         & 11,933         & 11,901         & 11,382\\  
      R$^2$                     & 0.94768        & 0.93700        & 0.92620        & 0.92320\\  
      Adjusted R$^2$            & 0.94537        & 0.93426        & 0.92297        & 0.91969\\  
      \bottomrule
   \end{tabular}
   \par \raggedright 
   \caption{Estimation for submitted bids under different outlier elimination rules. This table shows estimates for submitted bids according to  regression \eqref{eq:reg_sub_bid} for auctions that awarded at least one bid. Each model is defined by the type of method used for outlier detection according to section \ref{app:data}. Standardized product and region fixed effect are included in all Models. We compute \textit{robust} standard errors.}
\end{table}




\section{Regressions for awarded bids}


\begin{table}[H]
   \centering
   \begin{tabular}{lcccc}
      \toprule
                                & Type 1         & Type 2         & Type 3         & Type 4 \\   
                                & (1)            & (2)            & (3)            & (4)\\  
      \midrule 
      New                       & -0.072$^{***}$ & -0.073$^{***}$ & -0.064$^{***}$ & -0.055$^{***}$\\   
                                & (0.005)        & (0.005)        & (0.006)        & (0.007)\\   
      New $\times$ Comp  & -0.162$^{***}$ & -0.173$^{***}$ & -0.074$^{***}$ & -0.081$^{***}$\\   
                                & (0.008)        & (0.009)        & (0.012)        & (0.013)\\   
       \\
      Observations              & 11,760         & 11,933         & 11,901         & 11,382\\  
      R$^2$                     & 0.94555        & 0.93189        & 0.89614        & 0.89202\\  
      Adjusted R$^2$            & 0.94314        & 0.92893        & 0.89160        & 0.88708\\  
      \bottomrule
   \end{tabular}
   
   \par \raggedright 
   \caption{Estimation for awarded bids  under different outlier elimination rules. This table shows estimates for awarded bids. Each model is defined by the type of method used for outlier detection according to Appendix \ref{app:data}. Standardized product and region fixed effect are included in all models. We compute \textit{robust} standard errors.}
\end{table}




\begin{table}[H]
   \label{app:tab:volume_adj_awarded}
   \centering
   \begin{tabular}{lcccc}
      \toprule
                                & Type 1         & Type 2         & Type 3         & Type 4 \\   
                                & (1)            & (2)            & (3)            & (4)\\  
      \midrule 
      New                       & -0.033$^{***}$ & -0.035$^{***}$ & -0.105$^{***}$ & -0.105$^{***}$\\   
                                & (0.005)        & (0.005)        & (0.006)        & (0.007)\\   
      New $\times$ Comp  & -0.156$^{***}$ & -0.167$^{***}$ & -0.082$^{***}$ & -0.092$^{***}$\\   
                                & (0.008)        & (0.009)        & (0.011)        & (0.012)\\   
       \\
      Observations              & 11,760         & 11,933         & 11,901         & 11,382\\  
      R$^2$                     & 0.94506        & 0.93146        & 0.89717        & 0.89362\\  
      Adjusted R$^2$            & 0.94263        & 0.92848        & 0.89268        & 0.88875\\  
      \bottomrule
   \end{tabular}
   
   \par \raggedright 
   \caption{\label{tab:estadj_vol} Estimation for awarded bids: Adjusting by volume. This table shows estimates for awarded bids when prices are adjusted by volume following \eqref{eq:reg_sub_bid_rob}. Each model is defined by the type of method used for outlier detection according to Appendix~\ref{app:data}. Standardized product and region fixed effect are included in all models. We compute \textit{robust} standard errors.}
\end{table}




\section{Regressions for posted prices}


\begin{table}[H]
   \centering
   \begin{tabular}{lccc}
      \toprule
                                & Type 1         & Type 2         & Type 3 \\   
                                & (1)            & (2)            & (3)\\  
      \midrule 
      New                       & -0.029$^{***}$ & -0.033$^{***}$ & -0.025$^{***}$\\   
                                & (0.0007)       & (0.0008)       & (0.0008)\\   
      New $\times$ Comp  & -0.076$^{***}$ & -0.075$^{***}$ & -0.092$^{***}$\\   
                                & (0.0009)       & (0.0009)       & (0.001)\\   
       \\
      Observations              & 1,012,598      & 1,018,778      & 973,195\\  
      R$^2$                     & 0.97025        & 0.96865        & 0.96132\\  
      Adjusted R$^2$            & 0.97018        & 0.96857        & 0.96121\\  
      \bottomrule
   \end{tabular}
   
   \par \raggedright 
   \caption{Estimation for posted prices  under different outlier elimination rules. This table shows estimates for posted prices as in \eqref{eq:reg_op_posted}. Each Model is defined by the type of method used for outlier detection according to Appendix~\ref{app:data}. Standardized product, region, and week category-specific fixed effect are included in all models. We compute \textit{robust} standard errors.}
\end{table}



\section{Regressions for transaction prices}


\begin{table}[H]
   \centering
   \begin{tabular}{lccc}
      \toprule
                                & Type 1         & Type 2         & Type 3 \\   
                                & (1)            & (2)            & (3)\\  
      \midrule 
      New                       & -0.007$^{***}$ & -0.006$^{***}$ & 0.004$^{**}$\\   
                                & (0.002)        & (0.002)        & (0.002)\\   
      New $\times$ Comp  & -0.068$^{***}$ & -0.071$^{***}$ & -0.082$^{***}$\\   
                                & (0.002)        & (0.002)        & (0.002)\\   
       \\
      Observations              & 193,702        & 194,312        & 180,421\\  
      R$^2$                     & 0.97441        & 0.97328        & 0.97395\\  
      Adjusted R$^2$            & 0.97409        & 0.97294        & 0.97360\\  
      \bottomrule
   \end{tabular}
   
   \par \raggedright 
   \caption{Estimation for transaction prices. This table shows estimates for transaction prices as defined in Equation \eqref{eq:reg_op_posted}. Each model is defined by the type of method used for outlier detection according to Appendix~\ref{app:data}. Standardized product, region and week category-specific fixed effect are included in all models. We compute \textit{robust} standard errors.}
\end{table}



 % opcionales

\end{document}
